\documentclass{article}
\usepackage{graphicx} % Required for inserting images
\usepackage{amsmath,amsthm,amssymb,amsxtra} %Required for mathematical equations and symbols
\usepackage{mathtools} %Add-on to amsmath

\title{Solution to American Mathematical Monthly Problem 12598}
\author{Jayden Brown-Akoto}
\date{August 2026}

\begin{document}



\begin{figure}
    \centering
    \includegraphics[width=0.8\linewidth]{Screenshot 2026-08-16 2.29.22 PM.png}
    \label{fig:placeholder}
\end{figure}
\subsubsection*{Solution to Problem 12598 in Amer. Math. Monthly 133 (2026)}
\begin{center}
    Jayden Brown-Akoto (student author)
\end{center}

We show that the desired sums are $\ln(2)/4+\pi/8$ and $\pi^{2}/32+\pi/8-\ln(2)/4$, respectively.

\vspace{5mm}

\textbf{(a)} Let 

\begin{equation*}
    S := \sum_{n=1}^{\infty}(-1)^{n-1}\left(\sum_{k=0}^{\infty}\frac{(-1)^k}{n+2k}\right).
\end{equation*}

Then, we can rewrite $S$ to get

\begin{equation*}
    S =\sum_{n=1}^{\infty}\sum_{k=0}^{\infty}\frac{(-1)^{n-1}(-1)^k}{n+2k} =\sum_{n=1}^{\infty}\sum_{k=0}^{\infty}\int_{0}^{\infty}(-1)^{n-1}(-1)^ke^{-(n+2k)t}dt.
\end{equation*}

\vspace{5mm} 

To justify the interchange of integration and double summation, we apply the dominated convergence theorem twice. For the first swap, we define the sequence of functions

\begin{align*}
    f_{n, K}(t):=\sum_{k=0}^{K}(-1)^{n-1}(-1)^ke^{-(n+2k)t}dt.
\end{align*}

\vspace{5mm} 

We note that as $K \to \infty$ the sequence of functions converges pointwise for all $t>0$. 

\vspace{5mm} 

Then
\begin{align*}
    |f_{n, K}(t)|&=\left|\sum_{k=0}^{K}(-1)^{n-1}(-1)^ke^{-(n+2k)t}\right| = \left|e^{-nt}\cdot\frac{1-(-e^{-2t})^{(K+1)}}{1+e^{-2t}}\right| \\ &\leq \left|e^{-nt}\cdot\frac{1+e^{-2t(K+1)}}{1+e^{-2t}}\right|= e^{-nt}\cdot\frac{1+e^{-2t(K+1)}}{1+e^{-2t}} \leq e^{-nt}
\end{align*}

 for all $n\geq1$, $K \geq 0$, and $t\geq0$. Additionally, $\int_0^{\infty}e^{-nt}dt=1/n$ for all $n\geq1$. Thus, by the dominated convergence theorem, 
\begin{equation*}
    S=\sum_{n=1}^{\infty}\int_{0}^{\infty}\sum_{k=0}^{\infty}(-1)^{n-1}(-1)^ke^{-(n+2k)t}dt.
\end{equation*}
 

\vspace{5mm}

To justify the second swap, we define a sequence of functions
\begin{align*}
    m_{N}(t):=\sum_{n=1}^{N}\sum_{k=0}^{\infty}(-1)^{n-1}(-1)^ke^{-(n+2k)t}dt.
\end{align*}

We note that as $N \to \infty$ the sequence converges pointwise for all $t>0$.

\vspace{5mm}

Then 
\begin{align*}
    |m_N(t)|&=\left|\sum_{n=1}^{N}\sum_{k=0}^{\infty}(-1)^{n-1}(-1)^ke^{-(n+2k)t}\right| = \left|e^{-t}\cdot\frac{1-(-e^{-t})^N}{(1+e^{-t})(1+e^{-2t})}\right|\\ &\leq \left|e^{-t}\cdot   \frac{1+e^{-Nt}}{(1+e^{-t})(1+e^{-2t})}\right| = e^{-t}\cdot   \frac{1+e^{-Nt}}{(1+e^{-t})(1+e^{-2t})} \leq  e^{-t}
\end{align*}

for all $N \geq 1$ and $t\geq0$. Additionally, $\int_0^{\infty}e^{-t}dt=1$. Thus, by the dominated convergence theorem, 
\begin{equation*}
    S=\int_{0}^{\infty}\sum_{n=1}^{\infty}\sum_{k=0}^{\infty}(-1)^{n-1}(-1)^ke^{-(n+2k)t}dt.
\end{equation*}

\vspace{5mm}

Evaluating the double sum, we get

\begin{align*}
    S &= \int_{0}^{\infty}\sum_{n=1}^{\infty}\sum_{k=0}^{\infty}(-1)^{n-1}(-1)^ke^{-(n+2k)t}dt \\
    &= \int_{0}^{\infty}\left(\sum_{n=1}^{\infty}(-1)^{n-1}e^{-nt}\right)\left(\sum_{k=0}^{\infty}(-1)^ke^{-2kt}\right)dt \\
    &=\int_{0}^{\infty}\left(\frac{e^{-t}}{1+e^{-t}}\right)\left(\frac{1}{1+e^{-2t}}\right)dt \\
    &= \int_{0}^{\infty}\frac{e^{-t}}{(1+e^{-t})(1+e^{-2t})}dt.
\end{align*}

\vspace{5mm}

Now, we rearrange the integrand and perform the substitution $x=1+e^{-t}$ to obtain

\begin{align*}
    \int_{0}^{\infty}\frac{e^{-t}}{(1+e^{-t})(1+e^{-2t})}dt &= -\int_{0}^{\infty}\frac{-e^{-t}}{(1+e^{-t})((1+e^{-t})(-1+e^{-t})+2)}dt \\ 
    &=-\int_{2}^{1}\frac{1}{x(x(x-2)+2)}dx \\
    &= \int_{1}^{2}\frac{1}{x(x^2-2x+2)}dx.
\end{align*}

We then decompose the integrand to obtain 

\begin{align*}
     \int_{1}^{2}\frac{1}{x(x^2-2x+2)}d&x = \int_{1}^{2}\left(\frac{1}{2x}+\frac{1}{x^2-2x+2}-\frac{x}{2(x^2-2x+2)}\right)dx \\
     &=\frac{1}{2}\int_{1}^{2}\frac{1}{x}dx+\int_{1}^{2}\frac{1}{x^2-2x+2}dx-\frac{1}{2}\int_{1}^{2}\frac{x}{x^2-2x+2}dx.
\end{align*}

\vspace{5mm}

We denote the result as $I_1+I_2+I_3$.

\vspace{5mm}

Evaluating $I_1$, we get
\begin{align*}
    I_1=\frac{1}{2}\int_{1}^{2}\frac{1}{x}dx=\frac{1}{2}(\ln(2)-\ln(1))=\frac{\ln(2)}{2}.
\end{align*}

Evaluating $I_2$, we get 
\begin{align*}
    I_2&=\int_{1}^{2}\frac{1}{x^2-2x+2}dx = \int_{1}^{2}\frac{1}{(x-1)^2+1}dx \\
     &=\arctan(2-1)-\arctan(1-1) = \arctan(1)=\frac{\pi}{4}.
\end{align*}

To evaluate $I_3$, we begin by integrating by parts with $u=x$ and \newline $v'={1}/({x^2-2x+2})$. Hence, 

\begin{align*}
    I_3=-\frac{1}{2}\int_{1}^{2}\frac{x}{x^2-2x+2}dx&=-\frac{x\arctan(x-1)}{2}\Big|_{1}^{2}+\frac{1}{2}\int_{1}^{2}\arctan(x-1)dx \\
     &= -\frac{\pi}{4}+\frac{1}{2}\int_{1}^{2}\arctan(x-1)dx.
\end{align*}

For the remaining integral, we make the substitution $x=\tan(\theta)+1$ and obtain
\begin{align*}
   \frac{1}{2}\int_{1}^{2}\arctan(x-1)dx=\frac{1}{2}\int_{0}^{\frac{\pi}{4}}\arctan(\tan(\theta))\sec^2(\theta)d\theta=\frac{1}{2}\int_{0}^{\frac{\pi}{4}}\theta\sec^2(\theta)d\theta.
\end{align*}

Then, we integrate by parts with $u=\theta$ and $v'=\sec^2(\theta)$ to get

\begin{align*}
   \frac{1}{2}\int_{0}^{\frac{\pi}{4}}\theta\sec^2(\theta)d\theta&=\frac{1}{2}\left(\theta\tan\theta\Big|_{0}^{\frac{\pi}{4}}-\int_{0}^{\frac{\pi}{4}}\tan\theta d\theta\right) \\ 
   &=\frac{1}{2}\left(\theta\tan\theta\Big|_{0}^{\frac{\pi}{4}}+\ln(|\cos\theta|)\Big|_{0}^{\frac{\pi}{4}}\right) \\
   &=\frac{1}{2}\left(\frac{\pi}{4}-\frac{\ln(2)}{2}\right) \\
   &= \frac{\pi}{8}-\frac{\ln(2)}{4}.
\end{align*}

Thus, $I_3=-\pi/4+\pi/8-\ln(2)/4$ and

\begin{align*}
   I_1+I_2+I_3=\frac{\ln(2)}{2}+\frac{\pi}{4}+\frac{\pi}{8}-\frac{\ln(2)}{4}-\frac{\pi}{4}=\frac{\ln(2)}{4}+\frac{\pi}{8}.
\end{align*}

Therefore,
\begin{align*}
   S=\sum_{n=1}^{\infty}(-1)^{n-1}\left(\sum_{k=0}^{\infty}\frac{(-1)^k}{n+2k}\right)=\frac{\ln(2)}{4}+\frac{\pi}{8}.
\end{align*}
\vspace{5mm}

\textbf{(b)} Let 

\begin{equation*}
    S := \sum_{n=1}^{\infty}(-1)^{n-1}\left(\sum_{k=0}^{\infty}\frac{(-1)^k}{n+2k}\right)^2
\end{equation*}

and 

\begin{equation*}
    J_n := \sum_{k=0}^{\infty}\frac{(-1)^k}{n+2k}.
\end{equation*}

Then 

\begin{align*}
    &S=\sum_{n=1}^{\infty}(-1)^{n-1}J_n^2.
\end{align*}

\vspace{5mm}

Additionally, $J_n$ can be rewritten as
\begin{align*}
    \sum_{k=0}^{\infty}\int_{0}^{\infty}(-1)^{k}e^{-(n+2k)s}ds.
\end{align*}

To justify swapping integration and summation, we take a similar approach as used to justify the first swap in part \textbf{(a)} and define a sequence of functions
\begin{align*}
    f_{n,K}(s):=\sum_{k=0}^{K}(-1)^ke^{-(n+2k)s}ds.
\end{align*}

We note that as $K \to \infty$ the sequence converges pointwise for all $s>0$.

\vspace{5mm}

Then 
\begin{align*}
    |f_{n, K}(s)|&=\left|\sum_{k=0}^{K}(-1)^ke^{-(n+2k)s}\right| = \left|e^{-ns}\cdot\frac{1-(-e^{-2s})^{(K+1)}}{1+e^{-2s}}\right| \\ &\leq e^{-ns}\cdot\left|\frac{1+e^{-2s(K+1)}}{1+e^{-2s}}\right| \leq e^{-ns}
\end{align*}
for all $K \geq 0$ and $s \geq 0$. Moreover, $\int_0^\infty e^{-ns}ds=1/n$. Thus, by the dominated convergence theorem,
\begin{align*}
    \sum_{k=0}^{\infty}\int_{0}^{\infty}(-1)^{k}e^{-(n+2k)s}ds=\int_{0}^{\infty}\sum_{k=0}^{\infty}(-1)^{k}e^{-(n+2k)s}ds.
\end{align*}

Hence,
\begin{align*}
    J_n^2=J_n\cdot J_n&=\left(\int_{0}^{\infty}\sum_{k=0}^{\infty}(-1)^{k}e^{-(n+2k)s}ds\right)\left(\int_{0}^{\infty}\sum_{k=0}^{\infty}(-1)^{k}e^{-(n+2k)t}dt\right) \\
    &= \int_{0}^{\infty}\int_{0}^{\infty}\left(\sum_{k=0}^{\infty}(-1)^{k}e^{-(n+2k)s}\right)\left(\sum_{k=0}^{\infty}(-1)^{k}e^{-(n+2k)t}\right)dsdt \\
    &= \int_{0}^{\infty}\int_{0}^{\infty}\left(\frac{e^{-ns}}{1+e^{-2s}}\right)\left(\frac{e^{-nt}}{1+e^{-2t}}\right)dsdt \\
    &=\int_{0}^{\infty}\int_{0}^{\infty}\frac{e^{-n(s+t)}}{(1+e^{-2s})(1+e^{-2t})}dsdt.
\end{align*}

\vspace{5mm}

Substituting $J_n^2$ into $S$, we get

\begin{align*}
    S &= \sum_{n=1}^{\infty}(-1)^{n-1}\int_{0}^{\infty}\int_{0}^{\infty}\frac{e^{-n(s+t)}}{(1+e^{-2s})(1+e^{-2t})}dsdt\\
     &= \sum_{n=1}^{\infty}\int_{0}^{\infty}\int_{0}^{\infty}\frac{(-1)^{n-1}e^{-n(s+t)}}{(1+e^{-2s})(1+e^{-2t})}dsdt.
\end{align*}

\vspace{5mm}

To justify swapping summation and double integration, we use the dominated convergence theorem again. 

\vspace{5mm}

We define a sequence of multivariate functions
\begin{align*}
    f_N(s, t)=\sum_{n=1}^{N}\frac{(-1)^{n-1}e^{-n(s+t)}}{(1+e^{-2s})(1+e^{-2t})}.
\end{align*}.

We note that as $N \to \infty$ the sequence converges pointwise for all $s,t>0$.

\vspace{5mm}

Then,

\begin{align*}
    \left|\sum_{n=1}^{N}\frac{(-1)^{n-1}e^{-n(s+t)}}{(1+e^{-2s})(1+e^{-2t})}\right| &\leq \left|e^{-(s+t)}\cdot\frac{1-\left(-e^{-(s+t)}\right)^N}{(1+e^{-(s+t)})(1+e^{-2s})(1+e^{-2t})}\right| \\
    &\leq \left|e^{-(s+t)}\cdot\frac{1+e^{-N(s+t)}}{(1+e^{-(s+t)})(1+e^{-2s})(1+e^{-2t})}\right| \\ &= e^{-(s+t)}\cdot\frac{1+e^{-N(s+t)}}{(1+e^{-(s+t)})(1+e^{-2s})(1+e^{-2t})} \\ &\leq e^{-(s+t)}
\end{align*}

for all $N \geq 1$ and $s,t \geq 0$. Additionally, $\int_{0}^{\infty}\int_{0}^{\infty}e^{-(s+t)}dsdt=1$. Thus, by the dominated convergence theorem,

\begin{align*}
    S = \int_{0}^{\infty}\int_{0}^{\infty}\sum_{n=1}^{\infty}\frac{(-1)^{n-1}e^{-n(s+t)}}{(1+e^{-2s})(1+e^{-2t})}dsdt.
\end{align*}
\vspace{5mm}

Evaluating the sum, we obtain

\begin{align*}
    S &= \int_{0}^{\infty}\int_{0}^{\infty}\sum_{n=1}^{\infty}\frac{(-1)^{n-1}e^{-n(s+t)}}{(1+e^{-2s})(1+e^{-2t})}dsdt\\
    &= \int_{0}^{\infty}\int_{0}^{\infty}\frac{e^{-(s+t)}}{(1+e^{-(s+t)})(1+e^{-2s})(1+e^{-2t})}dsdt.
\end{align*}

\vspace{5mm}

Now, we perform the transformation $s=-\ln u$, $t= -\ln v$. The Jacobian determinant $|\partial(s,t)/\partial(u,v)|$ of this transformation is 

\[
\begin{vmatrix}
-\frac{1}{u} & 0  \\
0 & -\frac{1}{v}
\end{vmatrix} = \frac{1}{uv}.
\]

\vspace{5mm}

Hence,

\begin{align*}
    S &= \int_{0}^{\infty}\int_{0}^{\infty}\frac{e^{-(s+t)}}{(1+e^{-(s+t)})(1+e^{-2s})(1+e^{-2t})}dsdt\\
    &= \int_{1}^{0}\int_{1}^{0}\frac{uv}{(1+uv)(1+u^2)(1+v^2)}\cdot\frac{1}{uv}dudv \\
    &= \int_{0}^{1}\int_{0}^{1}\frac{1}{(1+uv)(1+u^2)(1+v^2)}dudv. 
\end{align*}

\vspace{5mm}

Decomposing the integrand, we get
\begin{align*}
    S = &\int_{0}^{1}\int_{0}^{1}\left(\frac{v^2}{(1+uv)(1+v^2)^2}+\frac{1}{(1+u^2)(1+v^2)^2}-\frac{uv}{(1+u^2)(1+v^2)^2}\right)dudv \\
    = &\int_{0}^{1}\int_{0}^{1}\frac{1}{(1+u^2)(1+v^2)^2}dudv-\int_{0}^{1}\int_{0}^{1}\frac{uv}{(1+u^2)(1+v^2)^2}dudv \\&+\int_{0}^{1}\int_{0}^{1}\frac{v^2}{(1+uv)(1+v^2)^2}dudv.
\end{align*}

We denote the result as $I_1+-I_2+I_3$.

\vspace{5mm}

Simplifying $I_1$, we obtain
\begin{align*}
    I_1&=\int_{0}^{1}\int_{0}^{1}\frac{1}{(1+u^2)(1+v^2)^2}dudv \\&=\left(\int_{0}^{1}\frac{1}{1+u^2}du\right)\left(\int_{0}^{1}\frac{1}{(1+v^2)^2}dv\right) \\
    &=\arctan(1) \cdot \int_{0}^{1}\frac{1}{(1+v^2)^2}dv \\
    &=\frac{\pi}{4}\int_{0}^{1}\frac{1}{(1+v^2)^2}dv.
\end{align*}

\vspace{60mm}

Now, we perform the substitution $v=\tan\theta$ and evaluate the integral to get 

\begin{align*}
    \frac{\pi}{4}\int_{0}^{1}\frac{1}{(1+v^2)^2}dv&=\frac{\pi}{4}\int_{0}^{\frac{\pi}{4}}\frac{\sec^2\theta}{(1+\tan^2\theta)^2}d\theta \\
    &=\frac{\pi}{4}\int_{0}^{\frac{\pi}{4}}\frac{\sec^2\theta }{(\sec^2\theta)^2}d\theta \\
    &=\frac{\pi}{4}\int_{0}^{\frac{\pi}{4}}\frac{1}{\sec^2\theta}d\theta \\
    &=\frac{\pi}{4}\int_{0}^{\frac{\pi}{4}}\frac{1+\cos2\theta}{2}d\theta \\
    &=\frac{\pi}{4}\Big(\frac{2\theta+\sin2\theta}{4}\Big|_0^\frac{\pi}{4}\Big) \\
    &=\frac{\pi}{4}\left(\frac{\pi}{8}+\frac{1}{4}\right)=\frac{\pi^2}{32}+\frac{\pi}{16}.
\end{align*}

Thus, $I_1=\pi^2/32+\pi/16$.

\vspace{5mm}

Evaluating $-I_2$, we obtain
\begin{align*}
    -\int_{0}^{1}\int_{0}^{1}\frac{uv}{(1+u^2)(1+v^2)^2}dudv&=-\left(\int_{0}^{1}\frac{u}{1+u^2}du\right)\left(\int_{0}^{1}\frac{v}{(1+v^2)^2}dv\right) \\
    &=-\Big(\frac{1}{2}\ln(1+u^2)\Big|_0^1\Big) \cdot \Big(-\frac{1}{2(1+v^2)}\Big|_0^1\Big) \\
    &=-\frac{\ln(2)}{2}\cdot \frac{1}{4}=-\frac{\ln(2)}{8}.
\end{align*}

\vspace{5mm}

To evaluate $I_3$, we perform the transformation $u=(x-1)/y$, $v=y$. The Jacobian determinant $|\partial(u, v)/\partial(x,y)|$ of the transformation is 

\[
\begin{vmatrix}
\frac{1}{y} & -\frac{(x-1)}{y^2} \\
0 & 1
\end{vmatrix} = \frac{1}{y}.
\]

Therefore,
\begin{align*}
    I_3 = \int_{0}^{1}\int_{0}^{1}\frac{v^2}{(1+uv)(1+v^2)^2}dudv&=\int_{0}^{1}\int_{1}^{y+1}\frac{y}{x(1+y^2)^2}dxdy \\ &=\int_{0}^{1}\frac{y}{(1+y^2)^2}\left(\int_{1}^{y+1}\frac{1}{x}dx\right)dy.
\end{align*}

Evaluating the inner integral, we obtain
\begin{align*}
   \int_{0}^{1}\frac{y}{1+y^2}\left(\int_{1}^{y+1}\frac{1}{x}dx\right)dy=\int_{0}^{1}\frac{y\ln(y+1)}{(1+y^2)^2}dy.
\end{align*}

Then, we integrate by parts with $u=\ln(y+1)$ and $v'=y/(1+y^2)^2$ to obtain
\begin{align*}
   \int_{0}^{1}\frac{y\ln(y+1)}{(1+y^2)^2}dy&=-\frac{\ln(y+1)}{2(1+y^2)}\Big|_0^1+\frac{1}{2}\int_{0}^{1}\frac{1}{(1+y)(1+y^2)}dy \\
   &=-\frac{\ln(2)}{4}+\frac{1}{2}\int_{0}^{1}\frac{1}{(1+y)(1+y^2)}dy. 
\end{align*}
We then partially decompose the integrand of the remaining integral to obtain
\begin{align*}
  \frac{1}{2}\int_{0}^{1}\frac{1}{(1+y)(1+y^2)}&dy =\frac{1}{2}\int_0^1\left(\frac{1}{2(1+y)}-\frac{y}{2(1+y^2)}+\frac{1}{2(1+y^2)}\right)dy\\
   =&\frac{1}{4}\int_0^1\frac{1}{1+y}dy-\frac{1}{4}\int_0^1\frac{y}{1+y^2}dy+\frac{1}{4}\int_0^1\frac{1}{1+y^2}dy \\
   =&\frac{1}{4}\Big(\ln(1+y)\Big|_0^1\Big)-\frac{1}{4}\left(\frac{\ln(1+y^2)}{2}\Big|_0^1\right)+\frac{1}{4}\Big(\arctan(y)\Big|_0^1\Big) \\
   =&\frac{\ln(2)}{4}-\frac{\ln(2)}{8}+\frac{\pi}{16} \\
   =&\frac{\ln(2)}{8}+\frac{\pi}{16}.
\end{align*}
Thus, $I_3=-\ln(2)/4+\ln(2)/8+\pi/16$ and
\begin{align*}
  I_1+-I_2+I_3=\frac{\pi^2}{32}+\frac{\pi}{16}-\frac{\ln(2)}{8}-\frac{\ln(2)}{4}+\frac{\ln(2)}{8}+\frac{\pi}{16}=\frac{\pi^2}{32}+\frac{\pi}{8}-\frac{\ln(2)}{4}.
\end{align*}

\vspace{5mm}

Therefore,
\begin{align*}
  S=\sum_{n=1}^{\infty}(-1)^{n-1}\left(\sum_{k=0}^{\infty}\frac{(-1)^k}{n+2k}\right)^2=\frac{\pi^2}{32}+\frac{\pi}{8}-\frac{\ln(2)}{4}.
\end{align*}
\end{document}
`